\setcounter{section}{5}

  \zwischenueberschrift{Übungsaufgaben}

\inputaufgabegibtloesung
{}
{

Es sei

\mavergleichskettedisp
{\vergleichskette
{ f(x,y)
}
{ =} { ax^2+by^2+cxy
}
{ } {
}
{ } {
}
{ } {
}
}
{}{}{.}
Bestimme die
\definitionsverweis {Hauptkrümmungen}{}{}
und die
\definitionsverweis {Hauptkrümmungsrichtungen}{}{}
des
\definitionsverweis {Graphen}{}{}
zu $f$ im Nullpunkt.

}
{} {}

\inputaufgabe
{}
{

Bestimme die
\definitionsverweis {Hauptkrümmungen}{}{}
und die
\definitionsverweis {Hauptkrümmungsrichtungen}{}{}
für jeden Punkt des
\definitionsverweis {Graphen}{}{}
zu
\maabbeledisp {f} {\R^{n-1}} {\R
} { \left(  x_1   , \,   \ldots  , \,   x_{n-1}                 \right) } { a_1x_1^2  + \cdots + a_{n-1} x_{n-1}^2
} {.}

}
{} {}

\inputaufgabe
{}
{

Es sei

\mavergleichskettedisp
{\vergleichskette
{ f(x,y,z)
}
{ =} { ax^2 +by^2 +c z^2 +dxy
}
{ } {
}
{ } {
}
{ } {
}
}
{}{}{.}
Berechne das
\definitionsverweis {charakteristische Polynom}{}{}
der
\definitionsverweis {Weingartenabbildung}{}{}
des
\definitionsverweis {Graphen}{}{}
zu $f$ im Nullpunkt. Berechne ferner die
\definitionsverweis {Hauptkrümmungen}{}{}
und die
\definitionsverweis {Hauptkrümmungsrichtungen}{}{.}

}
{} {}

\inputaufgabe
{}
{

Es sei

\mavergleichskettedisp
{\vergleichskette
{ f(x,y,z)
}
{ =} { ax^2 +by^2 +c z^2 +dxy + e yz
}
{ } {
}
{ } {
}
{ } {
}
}
{}{}{.}
Berechne das
\definitionsverweis {charakteristische Polynom}{}{}
der
\definitionsverweis {Weingartenabbildung}{}{}
des
\definitionsverweis {Graphen}{}{}
zu $f$ im Nullpunkt. Wie schwierig ist es in diesem Fall, die
\definitionsverweis {Hauptkrümmungen}{}{}
zu bestimmen?

}
{} {}

\inputaufgabe
{}
{

Es sei

\mavergleichskettedisp
{\vergleichskette
{ f(x,y)
}
{ =} { x^2+y^3
}
{ } {
}
{ } {
}
{ } {
}
}
{}{}{.}
Bestimme die
\definitionsverweis {Hauptkrümmungen}{}{,}
die
\definitionsverweis {Hauptkrümmungsrichtungen}{}{}
und die
\definitionsverweis {Gaußkrümmung}{}{}
des
\definitionsverweis {Graphen}{}{}
zu $f$ für jeden Punkt

\mavergleichskette
{\vergleichskette
{ P
}
{ = }{ (Q,f(Q))
}
{  }{
}
{    }{
}
{   }{
}
}
{}{}{,}

\mavergleichskette
{\vergleichskette
{ Q
}
{ \in }{ \R^2
}
{  }{
}
{    }{
}
{   }{
}
}
{}{}{.}

}
{} {}

\inputaufgabe
{}
{

Es sei

\mavergleichskette
{\vergleichskette
{ V
}
{ \subseteq }{ \R^{n-1}
}
{  }{
}
{    }{
}
{   }{
}
}
{}{}{}
\definitionsverweis {offen}{}{}
und sei
\maabb {f} { V } {\R
} {}
eine zweimal
\definitionsverweis {stetig differenzierbare Funktion}{}{.}
Zeige, dass die
\definitionsverweis {Hauptkrümmungen}{}{}
und die
\definitionsverweis {Hauptkrümmungsrichtungen}{}{}
des
\definitionsverweis {Graphen}{}{}
zu $f$ in einem Punkt

\mavergleichskette
{\vergleichskette
{ P
}
{ = }{ (Q,f(Q))
}
{  }{
}
{    }{
}
{   }{
}
}
{}{}{}
nur von den ersten und den zweiten partiellen Ableitungen von $f$ in $Q$ abhängen.

}
{} {}

\inputaufgabe
{}
{

Es sei

\mavergleichskette
{\vergleichskette
{ W
}
{ \subseteq }{ \R^n
}
{  }{
}
{    }{
}
{   }{
}
}
{}{}{}
\definitionsverweis {offen}{}{,}
\maabb {h} { W } { \R
} {}
eine
\definitionsverweis {stetig differenzierbare Funktion}{}{}
und

\mavergleichskette
{\vergleichskette
{ Y
}
{ = }{ h^{-1}(c)
}
{  }{
}
{    }{
}
{   }{
}
}
{}{}{}
die
\definitionsverweis {Faser}{}{}
zu

\mavergleichskette
{\vergleichskette
{ c
}
{ \in }{ \R
}
{  }{
}
{    }{
}
{   }{
}
}
{}{}{,}
wobei $h$ in jedem Punkt von $Y$
\definitionsverweis {regulär}{}{}
sei. Wir fassen $h$ auch als eine Abbildung $\tilde{h}$  auf
\mathl{W \times \R^m}{} auf, wobei die hinteren $m$ Variablen nicht explizit in $\tilde{h}$ eingehen. Es sei

\mavergleichskettedisp
{\vergleichskette
{ Z
}
{ =} { \tilde{h}^{-1}(c)
}
{ =} { Y \times \R^m
}
{ \subseteq} { W \times \R^{m}
}
{ } {
}
}
{}{}{.}
Es sei

\mavergleichskette
{\vergleichskette
{ Q
}
{ \in }{ Z
}
{  }{
}
{    }{
}
{   }{
}
}
{}{}{}
ein Punkt über

\mavergleichskette
{\vergleichskette
{ P
}
{ \in }{ Y
}
{  }{
}
{    }{
}
{   }{
}
}
{}{}{.}
In welcher Beziehung stehen die
\definitionsverweis {Hauptkrümmungen}{}{}
und
\definitionsverweis {Hauptkrümmungsrichtungen}{}{}
von $Z$ in $Q$ und von $Y$ in $P$.

}
{} {}

Es sei

\mavergleichskette
{\vergleichskette
{ W
}
{ \subseteq }{ \R^n
}
{  }{
}
{    }{
}
{   }{
}
}
{}{}{}
\definitionsverweis {offen}{}{,}
\maabb {h} { W } { \R
} {}
eine zweifach
\definitionsverweis {stetig differenzierbare Funktion}{}{}
und

\mavergleichskette
{\vergleichskette
{ Y
}
{ = }{ h^{-1}(c)
}
{  }{
}
{    }{
}
{   }{
}
}
{}{}{}
die
\definitionsverweis {Faser}{}{}
zu

\mavergleichskette
{\vergleichskette
{ c
}
{ \in }{ \R
}
{  }{
}
{    }{
}
{   }{
}
}
{}{}{,}
wobei $h$ in jedem Punkt von $Y$
\definitionsverweis {regulär}{}{}
sei. Zu einem Punkt

\mavergleichskette
{\vergleichskette
{ P
}
{ \in }{ Y
}
{  }{
}
{    }{
}
{   }{
}
}
{}{}{}
nennt man die
\definitionsverweis {Determinante}{}{}
der
\definitionsverweis {Weingartenabbildung}{}{}
$L_P$ die
\definitionswort {Gauß-Kronecker-Krümmung}{}
von $Y$ in $P$.

Die Gauß-Kronecker-Krümmung ist eine direkte Verallgemeinerung der Gaußkrümmung.

\inputaufgabe
{}
{

Es sei

\mavergleichskette
{\vergleichskette
{ W
}
{ \subseteq }{ \R^n
}
{  }{
}
{    }{
}
{   }{
}
}
{}{}{}
\definitionsverweis {offen}{}{,}
\maabb {h} { W } { \R
} {}
eine zweifach
\definitionsverweis {stetig differenzierbare Funktion}{}{}
und

\mavergleichskette
{\vergleichskette
{ Y
}
{ = }{ h^{-1}(c)
}
{  }{
}
{    }{
}
{   }{
}
}
{}{}{}
die
\definitionsverweis {Faser}{}{}
zu

\mavergleichskette
{\vergleichskette
{ c
}
{ \in }{ \R
}
{  }{
}
{    }{
}
{   }{
}
}
{}{}{,}
wobei $h$ in jedem Punkt von $Y$
\definitionsverweis {regulär}{}{}
sei. Es sei

\mavergleichskette
{\vergleichskette
{ P
}
{ \in }{ Y
}
{  }{
}
{    }{
}
{   }{
}
}
{}{}{.}
Zeige, dass das Produkt der
\definitionsverweis {Hauptkrümmungen}{}{}
von $Y$ in $P$ gleich der
\definitionsverweis {Gauß-Kronecker-Krümmung}{}{}
von $Y$ in $P$ ist.

}
{} {}

\inputaufgabe
{}
{

Es sei

\mavergleichskettedisp
{\vergleichskette
{ h(x,y,z,w)
}
{ =} { x^2+y^3+z^4+w^5
}
{ } {
}
{ } {
}
{ } {
}
}
{}{}{}
und

\mavergleichskette
{\vergleichskette
{ Y
}
{ = }{ h^{-1}(1)
}
{  }{
}
{    }{
}
{   }{
}
}
{}{}{.}
Bestimme für den Punkt

\mavergleichskettedisp
{\vergleichskette
{ P
}
{ =} { \left(  1   , \,   -1  , \,   1 , \,   0                \right)
}
{ \in} { Y
}
{ } {
}
{ } {
}
}
{}{}{}
die
\definitionsverweis {Weingartenabbildung}{}{}
$L_P$, ihr
\definitionsverweis {charakteristisches Polynom}{}{}
und die
\definitionsverweis {Gauß-Kronecker-Krümmung}{}{}
von $Y$ in $P$.

}
{} {}

\inputaufgabe
{}
{

Es sei

\mavergleichskettedisp
{\vergleichskette
{ h(x,y,z)
}
{ =} { x^2+y^3+z^5
}
{ } {
}
{ } {
}
{ } {
}
}
{}{}{}
auf

\mavergleichskette
{\vergleichskette
{ W
}
{ = }{ \R^3 \setminus \{0\}
}
{  }{
}
{    }{
}
{   }{
}
}
{}{}{}
und

\mavergleichskette
{\vergleichskette
{ Y
}
{ = }{ h^{-1}(0)
}
{  }{
}
{    }{
}
{   }{
}
}
{}{}{.}
Bestimme die
\definitionsverweis {Normalkrümmung}{}{}
im Punkt

\mavergleichskette
{\vergleichskette
{ P
}
{ = }{ \begin{pmatrix}  1  \\ -1\\  0   \end{pmatrix}
}
{  }{
}
{    }{
}
{   }{
}
}
{}{}{}
in Richtung des normierten Tangentialvektors
\mathl{{ \frac{   1  }{  \sqrt{29}  } }  \begin{pmatrix}  3  \\ -2\\  4   \end{pmatrix}}{.}

}
{} {}

  \zwischenueberschrift{Aufgaben zum Abgeben}

\inputaufgabe
{4}
{

Es sei

\mavergleichskettedisp
{\vergleichskette
{ f(x,y)
}
{ =} { x^3 +x^2y+3xy^2-y^4
}
{ } {
}
{ } {
}
{ } {
}
}
{}{}{.}
Bestimme die
\definitionsverweis {Hauptkrümmungen}{}{,}
die
\definitionsverweis {Hauptkrümmungsrichtungen}{}{}
und die
\definitionsverweis {Gaußkrümmung}{}{}
des
\definitionsverweis {Graphen}{}{}
zu $f$ für jeden Punkt

\mavergleichskette
{\vergleichskette
{ P
}
{ = }{ (Q,f(Q))
}
{  }{
}
{    }{
}
{   }{
}
}
{}{}{,}

\mavergleichskette
{\vergleichskette
{ Q
}
{ \in }{ \R^2
}
{  }{
}
{    }{
}
{   }{
}
}
{}{}{.}

}
{} {}

\inputaufgabe
{4}
{

Es sei

\mavergleichskettedisp
{\vergleichskette
{ f(x,y,z)
}
{ =} { -7x^2 +3y^2 -5z^2+ 6xy
}
{ } {
}
{ } {
}
{ } {
}
}
{}{}{.}
Bestimme die
\definitionsverweis {Hauptkrümmungen}{}{}
und die
\definitionsverweis {Hauptkrümmungsrichtungen}{}{}
des
\definitionsverweis {Graphen}{}{}
zu $f$ im Nullpunkt.

}
{} {}

\inputaufgabe
{6}
{

Es sei

\mavergleichskettedisp
{\vergleichskette
{ h(x,y,z,w)
}
{ =} { x^2+yz+z^3+xyzw
}
{ } {
}
{ } {
}
{ } {
}
}
{}{}{}
und

\mavergleichskette
{\vergleichskette
{ Y
}
{ = }{ h^{-1}(2)
}
{  }{
}
{    }{
}
{   }{
}
}
{}{}{.}
Bestimme für den Punkt

\mavergleichskettedisp
{\vergleichskette
{ P
}
{ =} { \left(  1   , \,   1  , \,   1 , \,   -1                \right)
}
{ \in} { Y
}
{ } {
}
{ } {
}
}
{}{}{}
die
\definitionsverweis {Weingartenabbildung}{}{}
$L_P$, ihr
\definitionsverweis {charakteristisches Polynom}{}{}
und die
\definitionsverweis {Gauß-Kronecker-Krümmung}{}{}
von $Y$ in $P$.

}
{} {Tipp: Löse die Gleichung nach einer geeigneten Variablen auf und behandle die Hyperfläche als einen Graphen in den anderen Variablen.}

\inputaufgabe
{6 (2+2+2)}
{

Es sei

\mavergleichskettedisp
{\vergleichskette
{ h(x,y,z)
}
{ =} { x^3+y^3+z^3
}
{ } {
}
{ } {
}
{ } {
}
}
{}{}{}
auf

\mavergleichskette
{\vergleichskette
{ W
}
{ = }{ \R^3 \setminus \{0\}
}
{  }{
}
{    }{
}
{   }{
}
}
{}{}{}
und

\mavergleichskette
{\vergleichskette
{ Y
}
{ = }{ h^{-1}(0)
}
{  }{
}
{    }{
}
{   }{
}
}
{}{}{.}
Es sei

\mavergleichskette
{\vergleichskette
{ P
}
{ = }{ \begin{pmatrix}  1  \\ 1\\  - \sqrt[3]{2}   \end{pmatrix}
}
{ \in }{ Y
}
{    }{
}
{   }{
}
}
{}{}{}
und

\mavergleichskettedisp
{\vergleichskette
{ v
}
{ =} { \begin{pmatrix}  1  \\ -1\\  0   \end{pmatrix}
}
{ \in} { T_PY
}
{ } {
}
{ } {
}
}
{}{}{.}
\aufzaehlungdrei{Bestimme eine Ebenengleichung für die
\definitionsverweis {Normalebene}{}{}
zu $v$.
}{Bestimme die
\definitionsverweis {Normalkrümmung}{}{}
in $P$ in Richtung des normierten Tangentialvektors
\mathl{{ \frac{   v  }{  \Vert { v } \Vert  } }}{.}
}{Bestimme eine bogenparametrisierte Kurve
\maabbdisp {\gamma} { I } { Y
} {}
mit

\mavergleichskette
{\vergleichskette
{ \gamma (0)
}
{ = }{ P
}
{  }{
}
{    }{
}
{   }{
}
}
{}{}{,}

\mavergleichskette
{\vergleichskette
{ \gamma' (0)
}
{ = }{ { \frac{   v  }{  \Vert { v } \Vert  } }
}
{  }{
}
{    }{
}
{   }{
}
}
{}{}{.}
Bestätige
[[Hyperfläche/Normalkrümmung/Normalebene/Kurvenschnitt/Fakt|Satz 5.14]]
in dieser Situation.
}

}
{} {}

\inputaufgabe
{2}
{

Skizziere eine differenzierbare Fläche $Y$ im Raum und eine Normalebene $E$, deren Durchschnitt mit der Fläche eine Kurve ergibt, die nicht überall regulär ist.

}
{} {}

 [[Kategorie:Latexseite]]
